\section{Discusión}

Los resultados obtenidos permiten interpretar el impacto del Sistema de Carnetización Digital en la gestión administrativa de instituciones educativas. Se analiza la generalización de los hallazgos, costos/beneficios, comparación con literatura y implicaciones prácticas.

\subsection{Interpretación de Efectos}

El sistema demuestra efectividad en reducir ineficiencias manuales, alineándose con literatura (e.g., Hamzah et al., 2021; Hendrawan et al., 2025). La reducción del 70% en tiempos de registro valida QR estáticos como solución viable para contextos con recursos limitados. Sin embargo, vulnerabilidades de seguridad resaltan la necesidad de capas adicionales, como se propone en Nwabuwe et al. (2023).

La alta satisfacción de usuarios (85%) indica buena aceptación, pero dependencias de conectividad e infraestructura móvil limitan generalización a áreas rurales o con baja cobertura.

\subsection{Generalización y Validez}

Los hallazgos son generalizables a instituciones educativas similares, pero requieren validación en escalas mayores. La arquitectura modular facilita adaptación a otros sectores (empresas, eventos públicos), aunque seguridad debe fortalecerse. Validez interna se asegura por triangulación de datos (pruebas, literatura, métricas); externa mediante comparación con estudios similares.

\subsection{Costos y Beneficios}

**Beneficios**:
- Reducción de costos administrativos (menos papel, tiempo humano).
- Mejora en precisión y accesibilidad de datos.
- Escalabilidad sin inversiones masivas.

**Costos**:
- Desarrollo inicial (tiempo, recursos).
- Dependencia de dispositivos móviles y conectividad.
- Riesgos de seguridad no mitigados.

El balance es positivo para instituciones con procesos manuales ineficientes, pero requiere evaluación caso por caso.

\subsection{Comparación con Literatura}

Comparado con sistemas avanzados (e.g., blockchain en Dongre y Verma, 2021; IA en Zihan e Islam, 2024), el enfoque simple prioriza accesibilidad sobre robustez. Coincide con recomendaciones de Razzaq y Ghayyur (2023) para migraciones graduales. Diferencia de soluciones monolíticas (e.g., Matz et al., 2023) al integrar automatización desde el inicio.

\subsection{Implicaciones para Decisores y Equipos}

Para administradores educativos: El sistema ofrece transición digital asequible, pero planificar expansiones de seguridad. Para desarrolladores: Valida metodologías ágiles para proyectos con restricciones temporales. Futuras investigaciones deben explorar integraciones híbridas (QR + RFID) para mayor seguridad.

En resumen, el sistema representa un paso práctico hacia digitalización, con potencial de evolución basado en literatura revisada.